
\subsection{Schutzziele und Informationssicherheit}
Die IT-Sicherheit befasst sich mit dem Schutz von
Informationen sowie der Systeme, die diese Informationen
verarbeiten. Es wird ein Zustand angestrebt, bei dem das
betrachtete System die definierten Sicherheitsziele erfüllt
(\cite[16]{von_faber_it_2021}).
Als die drei wichtigsten Schutzziele wird meist die CIA-Triade
herangezogen (\cite[21]{liedtke_informationssicherheit_2022}).
Das Akronym CIA steht für „Confidentiality“ (Vertraulichkeit),
„Integrity“ (Integrität) und
„Availability“ (Verfügbarkeit) (\cite[12]{von_faber_it_2021}).

Eine zentrale Herausforderung besteht darin, dass auf jeder
Ebene des Systems
Schwachstellen existieren können, die zur Verletzung dieser
Schutzziele führen können.
(\cite[19]{liedtke_informationssicherheit_2022})

Zur Gewährleistung von Informationssicherheit, werden sowohl
organisatorische als auch technische Aspekte betrachtet. Der
Fokus liegt dabei häufig auf Software, sowie auf der
Applikations- und Betriebssystemebene.
Schwachstellen in darunter liegenden Bereichen,
beispielsweise bei physischen Zugriffen, können die Sicherheit
des gesamten Systems schwächen.

Um Hardware Ressourcen Effizient
zu nutzen, werden diese zwischen verschiedenen Aufgaben geteilt.
(\cite[44]{kamp_building_2004})
Entsprechend existieren verschiedenste Mechanismen, um den Zugriff
eines Programms auf die Ressourcen anderer Programme zu
verhindern und so die Schutzziele zu gewährleisten. Diese
Isolation kann in Software beispielsweise als \ac{VM} oder durch
das Betriebssystem realisiert werden, geschieht aber häufig in
enger Verknüpfung mit der Hardware.

\subsection{Side-Channel und \acfp{TEA}}
Mit der Entdeckung praktischer Cache-basierter Side-Channels,
wie beispielsweise in \cite{bernstein_cache-timing_2005},
wurde ein neuer Angriffsvektor sichtbar, der auf der
konzeptionellen Funktionsweise moderner Computersysteme
beruht.
Das Problem mit dieser Art von Angriff ist, dass Schwachstellen
in den Bereichen zwischen der Hard- und Software ausgenutzt
werden, die in verschiedensten
\ac{CPU} Architekturen und Implementierungen auftreten.
Mit der Entdeckung von \acp{TEA} durch \cite{kocher_spectre_2019}
und \cite{lipp_meltdown_2020} wurde dieser Angriffsvektor
weiter vergrößert. Da diese Angriffe grundlegende
Konzepte moderner \acp{CPU} ausnutzen, betrifft diese Art an
Schwachstellen sowohl x86 Intel und AMD als auch ARM
\acp{CPU}-Architekturen. (\cite[253,255]{canella_systematic_2019})
\acp{TEA} nutzen dieselben Eigenschaften wie
Cache-basierte Side-Channel für sogenannte Covert-Channel,
um die Informationen wiederzuerlangen. Obwohl außer \ac{CbCC}
potenziell auch contentionbasierte Covert-Channel
(\cite[2]{wang_covert_2006}) eingesetzt
werden könnten, wurden bei \acp{POC} von \acp{TEA}
bisher grundsätzlich \acp{CbCC} verwendet.
(\cite[4,6]{kocher_spectre_2019},
\cite[4]{lipp_meltdown_2020}, \cite[7]{schwarz_zombieload_2019},
\cite[6]{kim_slap_2025}, \cite[91]{van_schaik_ridl_2019}).

Da \acp{TEA} eine hohe Komplexität
und starke Hardware Abhängigkeit aufweisen, wird ihre praktische
Umsetzbarkeit außerhalb kontrollierter
Laborumgebungen häufig hinterfragt. Aus demselben Grund wird
das von diesen Angriffen ausgehende Risiko hinterfragt.
(\cite[293,296]{shepherd_transient_2022})
Besonders bei browserbasierten Angriffen, wie beispielsweise
\cite{kim_ileakage_2023}, \cite{kim_slap_2025} und
\cite{kim_flop_2025}, die eine relativ hohe Laufzeit aufweisen
und einen passiven Nutzer voraussetzen, besteht die Frage,
ob diese Angriffe praktisch von einem Angreifer genutzt werden
könnten. Dabei könnte derselbe Angriff ein deutlich höheres
Risiko darstellen, insofern seine Laufzeit erheblich reduziert
werden kann.

Da das grundlegende Problem ohne
deutliche Performanceeinbußen nicht verhindert werden kann,
existieren verschiedene Ansätze zur Entwicklung von Mitigationen
gegen Covert-Channels während Transient Execution. Dazu zählen
sowohl automatische statische Analysen von Programmen, mit denen
sich mögliche Gadgets für solche Angriffe identifizieren lassen
(\cite{brotzman_specsafe_2021}),
als auch Mitigationen gegen spezifische Covert-Channel während
Transient Execution
(\cite{cheng_speclfb_2024}).
Eine generelle Mitigation gegen Transient Execution und
Cache-basierte Side-Channel Angriffe ist jedoch praktisch nicht
möglich, ohne die ausgenutzten Optimierungen vollständig zu
entfernen. Die einzige allgemeine Möglichkeit besteht darin
Störungen durch Änderungen in den Cache Zuständen auszulösen,
um eine Übertragung deutlich zu erschweren oder unmöglich zu machen.
(\cite[6]{wang_covert_2006}, \cite[11]{su_survey_2021})
Da \acp{TEA} allerdings nicht ausschließlich auf
\ac{CbCC} angewiesen sind und Störungen allein nicht den Angriff
verhindern, werden häufig weitere Maßnahmen betrachtet, die
allerdings nur spezifische Angriffsszenarien abdecken.
Aus diesem Grund bleiben solche Angriffe ein Risiko,
welches nicht vollständig mitigiert werden kann. Daraus ergibt
sich die Frage, welche Angriffe unter realistischen Bedingungen
tatsächlich praktikabel sind.

\subsection{Zielstellung}
Ziel dieser Arbeit ist es, einen Überblick über
die technischen Möglichkeiten von \ac{CbCC} hinsichtlich ihrer
Grenzen im Kontext von Transient Execution zu erarbeiten.

Hierzu werden zwei Ansätze zur Implementierung von CbCC sowie
verschiedene bekannte Techniken in Kombination mit Transient
Execution grundlegend verglichen. Dabei sollen
Parameter identifiziert werden, die für die
Übertragungsgeschwindigkeit eines CbCC entscheidend sind. Auf
dieser Grundlage soll eine Einschätzung der Grenzen solcher
CbCC auf gegebener Hardware ermöglicht werden.
